\documentclass[10pt,a4paper]{amsart}
\usepackage[margin=19mm]{geometry}
\usepackage{amsmath,amssymb,mathtools}
\usepackage[scr=rsfs,cal=euler]{mathalfa}
\usepackage{xfrac}
\usepackage[unicode,hidelinks,bookmarks=false]{hyperref}
\newcommand{\ie}{i.e.\ }
\newcommand{\cf}{cf.\ }
\newcommand{\resp}{respectively\ }
\newcommand{\Subsection}{\subsection}
\providecommand{\nobreakdash}{\nobreak\hbox{-}}
\setlength{\emergencystretch}{2em}
\multlinegap0pt
\usepackage{amssymb}
\usepackage{mathtools}
\usepackage{float}
\newtheorem{proposition}{Proposition}
\newcommand{\Z}{\ensuremath{\mathbb Z}}
\newcommand{\bx}{\boldsymbol{x}}
\newcommand{\bzero}{\boldsymbol{0}}
\newcommand{\myprime}{\,\prime}
\title{Mutation--selection balance on an infinite trait space: confinement, drift and equilibrium}
\author{Phil. Pollett}
\date{Source: arXiv:2609.02086v1 (CC BY 4.0)}
\begin{document}
\maketitle

\noindent\emph{Source and licence.} Mutation--selection balance on an infinite trait space: confinement, drift and equilibrium. Version arXiv:2609.02086v1 (CC BY 4.0), distributed by arXiv under the Creative Commons Attribution 4.0 International licence, http://creativecommons.org/licenses/by/4.0/, as recorded in the arXiv licence metadata for this exact version and shown on the abstract page https://arxiv.org/abs/2609.02086v1. Full licence notice: \texttt{licenses/arxiv-2609.02086-CC-BY-4.0.txt}.\par\smallskip

\noindent\textit{Editorial scope.} Counted unit: the article's proposition notattained (source0.tex lines 1272--1280). The article's complete original proof of it is delivered with it (source0.tex lines 2548--2654, one \texttt{proof} environment of 107 lines, which the article places away from the statement). No mathematics is written, repaired or re-derived: the statement and the proof are the article's own lines, in the article's own notation, and no retained line is substituted or re-flowed.  Retained uncounted as the source-local context the proof needs: lines 1109--1112.  Nothing was omitted from any retained span, so the package carries no omission record.  No retained line contains a citation, so the wrapper carries no bibliography and no bibliography entry.  No retained line contains a figure environment, an image-inclusion command or drawing code, so no image is delivered and the package declares no asset dependency.  Editorial formatting only, in addition: an AMS \texttt{amsart} preamble that restates the article's own declarations and macros used by the retained lines, namely the environments \texttt{proposition} on separate counters, together with the standard packages those lines need.  The retained environments print in this document as Proposition 1 in order of appearance, whereas the article prints the same items with the numbering of its own full document.  The complete native source of the article as distributed in the arXiv e-print for this exact version is bundled unchanged as \texttt{sources/209172/source0.tex}.
\begin{equation}
\dot x_i=\lambda x_i(\rho_i-m),\qquad i\in\Z,
\label{eq:nomutation}
\end{equation} 

\begin{proposition}
\label{notattained}
Suppose that $p_0=1$, $\bx(0)\neq\bzero$, and
$\mu^*:=\inf\{\mu_i:x_i(0)>0\}$
is finite but is not attained on the initial support. If
$\lambda>\mu^*$, then $m(t)\to\rho^*:=1-\mu^*/\lambda$.
Furthermore, for every $i\in\Z$,
$\pi_i(t)=x_i(t)/m(t)\to0$, and hence $x_i(t)\to0$.
\end{proposition}

\begin{proof}[Proof of Proposition~\ref{notattained}]
For each $i$ such that $x_i(0)>0$, equation~\eqref{eq:nomutation},
viewed as a scalar linear equation, implies that $x_i(t)>0$ for all
$t\geq0$. Dividing by $x_i(t)$ and integrating from $0$ to $t$ gives
$$
x_i(t)=x_i(0)e^{-\mu_i t}B(t),
\qquad
B(t):=\exp\left(\int_0^t\lambda(1-m(s))\,ds\right).
$$
The same formula holds when $x_i(0)=0$, since in that case
$x_i(t)=0$ for all $t\geq0$. Let
$Z(t):=\sum_i x_i(0)e^{-\mu_i t}$, so that $m(t)=B(t)Z(t)$.
Since $B(t)>0$, set $Y(t):=B(t)^{-1}$. From
$B^{\myprime}(t)=\lambda(1-m(t))B(t)$ and $m(t)Y(t)=Z(t)$, we obtain
$Y^{\myprime}(t)=-\lambda Y(t)+\lambda Z(t)$, with $Y(0)=1$. Hence
$$
Y(t)=e^{-\lambda t}
\left(1+\lambda\int_0^t e^{\lambda s}Z(s)\,ds\right),
$$
and therefore
\begin{equation}
m(t)
=
\frac{e^{\lambda t}Z(t)}
{1+\lambda\int_0^t e^{\lambda s}Z(s)\,ds}.
\label{eq:massrepresentation}
\end{equation}
Let $\delta_i:=\mu_i-\mu^*$ and
$\beta:=\lambda-\mu^*>0$, and write
$Z(t)=e^{-\mu^*t}a(t)$, where
$a(t):=\sum_i x_i(0)e^{-\delta_i t}$. Equation
\eqref{eq:massrepresentation} then becomes
\begin{equation}
m(t)
=
\frac{e^{\beta t}a(t)}
{1+\lambda\int_0^t e^{\beta s}a(s)\,ds}.
\label{eq:massrepresentation2}
\end{equation}
We first show that $-a^{\myprime}(t)/a(t)\to 0$. For every $\varepsilon>0$, let
$c_\varepsilon:=\sum_{\{i:\delta_i\leq\varepsilon/2\}}x_i(0)$.
Since $\mu^*$ is the infimum of the mortality rates on the initial
support, $c_\varepsilon>0$, and hence
$a(t)\geq c_\varepsilon e^{-\varepsilon t/2}$.
We next justify differentiating the series defining $a(t)$. Fix
$t_0>0$. For $t\geq t_0$,
$x_i(0)e^{-\delta_i t}\leq x_i(0)$, while
$$
\delta_i x_i(0)e^{-\delta_i t}
\leq \frac{x_i(0)}{e t_0},
$$
because $\sup_{u\geq0}ue^{-ut_0}=1/(e t_0)$. Since
$\sum_i x_i(0)<\infty$, both the series defining $a$ and the series of
its derivatives converge uniformly on $[t_0,\infty)$. Termwise
differentiation is therefore justified, and
$$
-a^{\myprime}(t)=\sum_i\delta_i x_i(0)e^{-\delta_i t},
\qquad t>0.
$$
Split this sum according to whether $\delta_i\leq\varepsilon$ or
$\delta_i>\varepsilon$. The first part is at most $\varepsilon a(t)$.
For all sufficiently large $t$, the function
$u\mapsto ue^{-ut}$ is decreasing on $[\varepsilon,\infty)$.
Consequently, if $\delta_i>\varepsilon$, then
$\delta_i e^{-\delta_i t}\leq\varepsilon e^{-\varepsilon t}$, and
hence
$-a^{\myprime}(t) \leq \varepsilon a(t)+\varepsilon m(0)e^{-\varepsilon t}$.
Using $a(t)\geq c_\varepsilon e^{-\varepsilon t/2}$, we obtain
$$
0\leq-\frac{a^{\myprime}(t)}{a(t)}
\leq
\varepsilon+
\frac{\varepsilon m(0)}{c_\varepsilon}
e^{-\varepsilon t/2}
$$
for all sufficiently large $t$. Letting first $t\to\infty$ and then
$\varepsilon\downarrow0$ shows that $-a^{\myprime}(t)/a(t)\to 0$.

Now set $f(t):=e^{\beta t}a(t)$. Then
$f^{\myprime}(t)/f(t)=\beta+a^{\myprime}(t)/a(t)\to\beta$. Since $\beta>0$, we have
$f^{\myprime}(t)/f(t)\geq\beta/2$ for all sufficiently large $t$. It follows
that $f(t)\to\infty$. Moreover, $f^{\myprime}(t)>0$ eventually, so
l'H{\^o}pital's rule gives
$$
\frac{\int_0^t f(s)\,ds}{f(t)} \to \frac{1}{\beta}.
$$
Substitution into \eqref{eq:massrepresentation2} yields
$m(t)\to {\beta}/{\lambda} = 1-{\mu^*}/{\lambda} = \rho^*$.
Finally,
$\pi_i(t)=x_i(0)e^{-\delta_i t}/a(t)$. If $x_i(0)=0$, then
$\pi_i(t)=0$ for all $t\geq0$. Otherwise, $\delta_i>0$ because the
infimum is not attained on the initial support. Choose
$\varepsilon\in(0,\delta_i)$ and let
$d_\varepsilon:=\sum_{\{j:\delta_j\leq\varepsilon\}}x_j(0)$.
Again, the definition of $\mu^*$ implies that $d_\varepsilon>0$, and
$a(t)\geq d_\varepsilon e^{-\varepsilon t}$. Therefore,
$$
0\leq\pi_i(t)
\leq
\frac{x_i(0)}{d_\varepsilon}
e^{-(\delta_i-\varepsilon)t}
\to 0.
$$
Thus, $\pi_i(t)\to 0$ for every $i\in\Z$. Since
$x_i(t)=m(t)\pi_i(t)$ and $m(t)\to\rho^*$, it follows that
$x_i(t)\to 0$ for every $i\in\Z$.
\end{proof}

\end{document}
